\documentclass[journal]{IEEEtran}
\usepackage{cite}
\usepackage{amsmath,amssymb,amsfonts}
\usepackage{graphicx}
\usepackage{booktabs}
\usepackage{xcolor}
\usepackage{url}
\usepackage{hyperref}
\hypersetup{colorlinks=true,linkcolor=blue,citecolor=blue,urlcolor=blue}

\begin{document}

\title{Semantic Signal-Assisted Decision Support for\\Reverse Logistics Resource Recovery Under Uncertainty}

\author{[Author 1], \textit{Amazon}, Seattle, WA, USA \and [Author 2], [Affiliation], [City, Country]}

\maketitle

%---------------------------------------------------------------------------
\begin{abstract}
Disposition decisions in reverse logistics pipelines—whether to refurbish,
harvest components, recycle, or scrap a returned product—must be made under
substantial uncertainty about actual condition. Conventional practice resolves
this uncertainty through physical inspection alone, which costs \$50--200 per
asset and 30--60 minutes of technician time, creating a bottleneck that delays
recovery and misallocates scarce processing capacity.

This paper describes a decision-support framework that extracts semantic signals
from the unstructured text already accompanying returns—technician notes,
maintenance logs, customer descriptions—and incorporates them as informative
priors in a stochastic disposition model. The extraction layer is
extractor-agnostic: any classifier, from simple keyword patterns to large language
models, can supply the required interface without changes to the downstream
optimizer. An adaptive inspection policy then uses extractor confidence to
determine inspection depth per asset, skipping physical testing where text
evidence is strong and reserving full component-level inspection for genuinely
uncertain cases.

We validate the approach across three structurally distinct scenarios—enterprise
IT decommission, aviation maintenance, repair, and overhaul (MRO), and consumer
electronics returns—using a deliberately conservative deterministic keyword
extractor as a lower bound on semantic extraction quality. Across all three
scenarios, the framework improves or matches the recovery value of an inspect-
every-asset baseline (net TRV lift 2.9--13.7\%, paired Wilcoxon $p<0.001$,
30~seeds) while saving meaningful inspection cost. Keyword-extractor calibration
on the decoupled synthetic benchmark gives Pearson $r=0.47$ for descriptive IT
notes, $r=0.32$ for standardized aviation codes, and $r=0.36$ for colloquial
consumer text; an oracle-text upper bound on the same benchmark reaches
$r=0.73$, $r=0.54$, and $r=0.74$ respectively, confirming that framework value
scales with extractor quality as designed. All results reproduce from the
released Reverse Logistics Decision Benchmark (RLDB).
\end{abstract}

\begin{IEEEkeywords}
reverse logistics, semantic signal extraction, resource recovery, adaptive
inspection, circular economy, stochastic optimization, decision support
\end{IEEEkeywords}

%---------------------------------------------------------------------------
\section{Introduction}
\label{sec:intro}

Resource recovery through reverse logistics is a growing operational challenge.
The global IT asset disposition market reached \$11.4B in~2023~\cite{GlobalMktInsights},
yet disposition accuracy remains low: returned products are routinely routed to
suboptimal recovery paths because condition assessment depends on physical inspection
that cannot scale with volumes. Aviation MRO operators face similar constraints—
with annual MRO demand exceeding \$80B~\cite{OliverWyman}, even small reductions
in unnecessary inspection translate to significant savings. Consumer electronics
returns exceed 3.6 billion units annually~\cite{NRF}, mostly processed without
any condition-based routing.

The core problem is an information gap. Decision-makers must choose among recovery
options with limited visibility into actual product condition. Physical inspection
resolves this uncertainty but is expensive and slow. Meanwhile, the textual data
that accompanies most returns—technician notes, maintenance logs, customer-reported
symptoms—often contains condition-relevant signals that existing optimization
frameworks cannot consume because they require structured numerical inputs~\cite{Godichaud2022,Govindan2015}.

We close this gap by formalizing the integration of semantic signals into
stochastic disposition optimization. The key observation is that unstructured
text, systematically extracted and calibrated, functions as a noisy but
informative prior on product condition. This prior parameterizes Beta-distributed
yield models that the decision engine uses to rank and allocate assets under
capacity constraints.

\smallskip
\noindent\textbf{Validation approach.}
We evaluate across three scenarios chosen to maximize structural diversity:

\begin{itemize}
\item \textbf{S1 (IT Infrastructure):} Planned decommissions with descriptive
technician notes. Keyword extractor calibration reaches $r=0.47$ (decoupled
synthetic benchmark); an oracle-text reader achieves $r=0.73$, bounding the
ceiling for this text type. The framework produces the largest TRV lift (13.7\%)
and largest inspection-cost savings.

\item \textbf{S2 (Aviation MRO):} Safety-critical domain with FAA Part~145
inspection mandates and standardized maintenance codes. Text signals are weaker
($r=0.32$), and the framework operates in advisory mode with lower TRV lift
(2.9\%) and meaningful inspection savings (\$34K per batch).

\item \textbf{S3 (Consumer Electronics):} High-volume returns where per-unit
inspection is economically infeasible. Continuous yield prediction is limited
($r=0.36$), but categorical routing—sorting returns into quality tiers without
individual inspection—provides the primary value at scale.
\end{itemize}

\smallskip
\noindent\textbf{Contributions.} This paper makes three contributions:
\begin{enumerate}
\item A formal framework for integrating semantic signals from unstructured text
into stochastic reverse logistics disposition models, designed to be independent
of the extraction method.
\item An adaptive inspection policy that uses extractor confidence to dynamically
determine inspection depth, reducing inspection costs while maintaining recovery
quality.
\item A simulation study across three structurally diverse scenarios establishing
where semantic signals add value, where they serve an advisory role, and where
they enable otherwise infeasible triage. The RLDB benchmark is released for
full reproducibility.
\end{enumerate}

%---------------------------------------------------------------------------
\section{Background and Related Work}

\subsection{Reverse Logistics Decision Challenges}

Reverse logistics decisions are sequential and partially irreversible: once an
asset is shredded or a component extracted, the action cannot be undone.
Classical planning approaches include mixed-integer programming for disassembly
lot-sizing under lead time uncertainty~\cite{Godichaud2022} and stochastic
programming for closed-loop network design~\cite{Govindan2015}. These methods
assume structured numerical inputs specified before optimization. In practice,
much of the decision-relevant information resides in unstructured sources that
no existing framework can ingest without manual translation.

\subsection{Text-Based Condition Assessment}

Large language models have recently been applied to forward supply chain
operations. OptiGuide~\cite{Yang2023} translates natural language queries into
optimization code; InvAgent~\cite{Quan2024} coordinates multi-agent inventory
decisions. AlMahri et al.~\cite{AlMahri2024} apply zero-shot learning to supply
chain relationship extraction. Despite this body of work, no system applies
semantic text extraction to reverse logistics disposition under yield uncertainty.
Table~\ref{tab:positioning} summarizes where this work sits relative to
existing approaches.

\begin{table}[t]
\caption{Positioning Relative to Existing Approaches}
\label{tab:positioning}
\centering
\begin{tabular}{lcccc}
\toprule
Approach & Struct. & Unstruct. & Yield & Adapt. \\
         & Input   & Text      & Uncert. & Insp. \\
\midrule
MIP/Stochastic \cite{Godichaud2022,Govindan2015} & \checkmark & -- & \checkmark & -- \\
ML Prediction \cite{Li2023,Zhang2022} & \checkmark & -- & Impl. & -- \\
LLM for SC \cite{Yang2023,Quan2024} & \checkmark & \checkmark & -- & -- \\
Rule-based \cite{Thierry1995} & \checkmark & -- & -- & -- \\
This work & \checkmark & \checkmark & \checkmark & \checkmark \\
\bottomrule
\end{tabular}
\end{table}

%---------------------------------------------------------------------------
\section{Framework Design}

\subsection{Problem Setting}

Assets arrive in periodic batches for disposition. Each asset $m$ has structured
attributes (type, age, source site, BOM structure) and unstructured context
(inspection notes $\mathcal{T}(m)$, site history $H(s(m))$, active business
directives). The BOM is modeled as a three-level hierarchy: $L_0$ (whole unit),
$L_1$ (major assemblies), and $L_2$ (individual components). Disassembly from
$L_1$ to $L_2$ costs 3--5$\times$ more than $L_0$ to $L_1$ due to precision
required for socketed component extraction.

Two decisions are made per asset: $D_0$ (triage—accept for processing or scrap
as-is) and $D_2$ (disposition—disassembly depth and component allocation).
The expected recovery value for asset $m$ disassembled to level $l$ is:
\begin{equation}
V(m, l) = \sum_{c \in \mathrm{BOM}(m,l)} p_c \cdot \mathbb{E}[Y_c(m)] - C_l
\label{eq:value}
\end{equation}
where $p_c$ is the market price of component $c$, $C_l$ is the disassembly cost
to reach level $l$, and $Y_c(m) \sim \mathrm{Beta}(\alpha_c, \beta_c)$ is the
stochastic yield for component $c$ in asset $m$.

\subsection{System Architecture}

The framework has two layers. The \emph{Semantic Extraction Layer} converts
unstructured text into a context factor $\phi \in (0,1]$ and a confidence score
$\sigma \in [0,1]$. The \emph{Recovery Decision Engine} consumes these alongside
structured features to evaluate Eq.~(\ref{eq:value}) at each feasible disassembly
level via Monte Carlo sampling ($N=1000$ draws), rank assets by expected net
value, and allocate greedily subject to capacity constraints.

\begin{figure}[t]
\centering
\fbox{\parbox{0.9\columnwidth}{\centering\small
\textbf{Extraction Layer} \\[2pt]
Unstructured text $\longrightarrow$ $\{\phi, \sigma\}$ \\[6pt]
$\downarrow$ \\[6pt]
\textbf{Adaptive Inspection Policy} ($\sigma$ vs $\tau_h, \tau_l$) \\[6pt]
$\downarrow$ \\[6pt]
\textbf{Recovery Decision Engine} (greedy + MC) \\[6pt]
$\downarrow$ \\[6pt]
Disposition + rationale
}}
\caption{Two-layer framework. The extraction layer maps text to $\phi$ (condition
factor) and $\sigma$ (confidence). The adaptive policy uses $\sigma$ to decide
inspection depth. The decision engine optimizes disposition under capacity
constraints using the adjusted yield prior.}
\label{fig:arch}
\end{figure}

\subsection{Context Extraction Layer}

The framework defines a semantic interface that decouples text extraction from
decision optimization. Any extractor satisfying this interface—returning
$\phi \in (0,1]$ and $\sigma \in [0,1]$—can serve as a drop-in component
without modifying the inspection policy or decision engine.

\smallskip
\noindent\textbf{Experimental extractor.} For reproducibility and to establish
a lower bound on framework performance, we implement a deterministic
keyword-and-pattern classifier with vocabularies drawn from iFixit teardown
reports (S1), ATA-coded SDR language (S2), and consumer return phrasing (S3).
The classifier is intentionally conservative: it falls back to the uninformative
prior ($\phi=1.0$) on any note for which no vocabulary signal fires. Calibration
on the decoupled synthetic benchmark (Section~\ref{sec:calibration}) gives
$r=0.47$ on S1, $r=0.32$ on S2, and $r=0.36$ on S3. An oracle-text reference
reader on the same benchmark achieves $r=0.73/0.54/0.74$, establishing the
ceiling reachable by any deterministic classifier on these text types.

\smallskip
\noindent\textbf{Why a deterministic extractor.} We choose a keyword classifier
for three reasons: (1)~full reproducibility without API dependency or model
versioning; (2)~clean attribution of performance gains to the adaptive policy
rather than extraction quality; and (3)~the known failure modes of pattern
matching—inability to parse negation, implicit context, domain abbreviations—
stress-test the downstream policy's robustness to noisy inputs. LLM-based
extractors would improve calibration; the interface requires no modification.

The context factor $\phi(m)$ adjusts baseline yield downward when negative
signals are present. Positive signals do not boost yield above baseline
($\phi \leq 1.0$). Multiple negative signals combine multiplicatively. When
extraction fails, the mechanism defaults to $\phi=1.0$, triggering full
inspection via the adaptive policy.

\subsection{Prior-Adjustment Mechanism}
\label{sec:framework-prior}

Given extractor output $\phi$ and baseline yield $\hat{y}_c$ for component $c$,
the Beta prior is adjusted by moment matching:
\begin{equation}
\alpha'_c = \phi(m) \cdot \hat{y}_c \cdot \kappa, \quad
\beta'_c  = (1 - \phi(m) \cdot \hat{y}_c) \cdot \kappa
\label{eq:prior}
\end{equation}
where $\kappa$ controls the effective prior sample size. This shifts the prior
mean toward the extractor's point estimate with strength proportional to source
reliability, without performing a formal Bayesian update—there is no likelihood
function or posterior computation. The mechanism is analogous to setting a
pseudo-count that encodes text-derived belief.

\smallskip
\noindent\textbf{Choice of $\kappa$.}
We set $\kappa=10$ as a practical default. Sensitivity analysis sweeping
$\kappa \in \{5, 10, 20, 40, 60, 80\}$ on S1 shows TRV lift ranging from
13.70\% to 14.56\%—a spread of under 1 percentage point
(Table~\ref{tab:kappa}). This near-flatness is analytically expected: $\kappa$
scales $\alpha'_c$ and $\beta'_c$ but leaves the Beta \emph{mean}
($\phi \cdot \hat{y}_c$) invariant, and the decision engine ranks assets by the
Monte Carlo mean. Framework value therefore derives from \emph{which} assets are
inspected, not from how tightly the prior concentrates around the text estimate.

\subsection{Adaptive Inspection Policy}

Rather than applying uniform inspection to every asset, the system uses
extractor confidence $\sigma(m)$ to determine inspection depth:
\begin{itemize}
\item $\sigma \geq \tau_h$: skip inspection, use $\phi$ directly;
\item $\tau_l \leq \sigma < \tau_h$: quick functional test;
\item $\sigma < \tau_l$: full component-level inspection.
\end{itemize}
Empirically, $\tau_h=0.5$ and $\tau_l=0.25$ work well for high-value assets
(S1, S2); S3 uses a narrower gap ($\tau_l=0.45$) because colloquial text
rarely provides very high confidence. When the extractor is confident, each
asset saves \$50--200 in inspection cost and 30--60 minutes of technician time.

\subsection{Decision Engine}

The Recovery Decision Engine processes weekly batches (daily for S3). For each
asset, it evaluates $V(m,l)$ at each feasible disassembly level via Monte Carlo
sampling of yield outcomes. Assets are ranked by expected net value and allocated
greedily to demand streams subject to weekly processing capacity. Under tight
capacity, the engine processes only the highest-value assets—a selective
harvesting strategy that emerges naturally from the capacity constraint.

%---------------------------------------------------------------------------
\section{Cross-Industry Evaluation}

\subsection{Scenario Design}

Table~\ref{tab:scenarios} describes the three scenarios. Each represents a
distinct class of reverse logistics operation: planned enterprise decommission,
condition-driven MRO, and volume-driven consumer returns. The same framework
architecture applies to all three; only the extraction vocabulary and economic
parameters differ.

\begin{table}[t]
\caption{Evaluation Scenarios}
\label{tab:scenarios}
\centering
\begin{tabular}{lccc}
\toprule
& S1: IT Infra. & S2: Aviation & S3: Consumer \\
\midrule
SC type & Product & Service & Product \\
RL driver & Planned decom & Oper.\ degrad. & Returns \\
Core uncert. & Comp.\ yield & Cond.\,+\,RUL & Prod.\ cond. \\
Text type & Tech.\ jargon & Maint.\ codes & Nat.\ lang. \\
Insp.\ cost & \$75, 60\,min & \$100, 35\,min & \$20, 5\,min \\
Assets & 500/batch & 500/batch & 1000/batch \\
Regulation & Minimal & FAA Pt.\ 145 & CPSC \\
\bottomrule
\end{tabular}
\end{table}

\subsection{Data Sources and Generation}

All simulation data is synthetic, calibrated from public sources
(Table~\ref{tab:datasources}). Asset prices, base yield rates, disassembly
costs, and note distributions are derived from public databases; the simulation
output is not tuned to hit any target.

A key design property is \emph{decoupled note generation}: in all three
scenarios, the technician note or review text is generated from an independently
drawn, noisy observation of ground-truth condition rather than from the condition
that sets true yield. Two corruption modes apply—omission (the note becomes
uninformative, probability 0.15) and severity mislabeling (perceived severity
drifts by one adjacent category, probability 0.25). This ensures that any
extractor--yield correlation reflects genuine signal recovery rather than a
construction artifact. See \texttt{src/data\_generators/noise.py} in the
released benchmark.

\begin{table}[t]
\caption{Parameter Calibration Sources}
\label{tab:datasources}
\centering\small
\begin{tabular}{lll}
\toprule
Parameter & Source & Scenario \\
\midrule
Component prices & PCPartPicker, eBay completed & S1 \\
Yield by age & Backblaze Drive Stats~\cite{Backblaze} & S1 \\
\quad{}          & Eurostat WEEE stats & S1 \\
Disassembly costs & BLS NAICS 562920, iFixit & S1 \\
SDR condition dist. & FAA SDR database~\cite{FAASDR} & S2 \\
Part prices & ILS, PartBase, HEICO & S2 \\
MRO costs & FAA Part~145 rate surveys & S2 \\
Return condition & CPSC SaferProducts~\cite{CPSC} & S3 \\
Refurb values & Back Market, Decluttr & S3 \\
Note vocabulary & iFixit teardowns, ATA codes & S1, S2 \\
\quad{}           & Amazon review corpus~\cite{Hou2024} & S3 \\
\bottomrule
\end{tabular}
\end{table}

\subsection{Baselines and Metrics}

We compare against five baselines: (1)~\emph{Random disposition}—random
assignment of $L_2$/$L_1$/scrap without ranking (lower bound); (2)~\emph{Rule-based
FIFO}—whole-unit processing in arrival order, no disassembly optimization
(industry-standard approximation); (3)~\emph{XGBoost proxy}—greedy optimizer
with $\phi$ derived from age bracket alone, learned on a held-out training
population (seed 99999, $N=2000$ assets) independent of the evaluation seeds,
no text; (4)~\emph{Optimization-only} ($\phi=1.0$, inspect all)—the strongest
inspect-everything baseline; and (5)~\emph{Semantic-only}—text signals without
the optimizer, threshold-based routing on $\phi$.

Three primary metrics: \emph{Total Recovery Value} (TRV), net dollars recovered
after all processing, inspection, and disposal costs; \emph{Demand Coverage Rate}
(DCR), fraction of incoming assets routed to a recovery pathway; and
\emph{Inspection Cost Savings} (ICS), dollars saved by the adaptive policy
relative to inspect-everything.

\subsection{Results}
\label{sec:results}

Table~\ref{tab:main} reports the main results. All figures are means over
30~independent seeds; each seed controls batch composition, per-asset yield
draws, and demand realization via a single master RNG with independent child
streams (no global state). The same seed is used across all methods, so
paired comparison is exact.

\begin{table*}[t]
\caption{Net Recovery Performance Across Scenarios (USD K; mean over 30 seeds)}
\label{tab:main}
\centering
\begin{tabular}{lrrrrrrrrr}
\toprule
& \multicolumn{3}{c}{S1: IT Infrastructure} &
  \multicolumn{3}{c}{S2: Aviation MRO} &
  \multicolumn{3}{c}{S3: Consumer Electronics} \\
\cmidrule(lr){2-4}\cmidrule(lr){5-7}\cmidrule(lr){8-10}
System & TRV & DCR & ICS & TRV & DCR & ICS & TRV & DCR & ICS \\
\midrule
Random        &  299 & 67\% &  13 & 1,077 & 66\% & 17 & $-$12 & 66\% &  7 \\
Rule-based    &    7 &100\% &   0 &    85 &100\% &  0 & $-$25 &100\% &  0 \\
XGBoost       &  570 & 78\% &   8 & 1,678 & 52\% & 24 &    71 & 87\% &  1 \\
Opt-only      &  533 & 78\% &   8 & 1,565 & 52\% & 24 &    68 & 96\% &  1 \\
Sem.-only     &  551 & 95\% &  32 & 1,377 & 75\% & 34 & $-$14 & 65\% &  5 \\
\textbf{Ours} &\textbf{606}&\textbf{100\%}&\textbf{32}&\textbf{1,610}&\textbf{57\%}&\textbf{34}&\textbf{71}&\textbf{94\%}&\textbf{1} \\
\midrule
Lift & \multicolumn{3}{c}{$+73$ ($+13.7\%$)} &
       \multicolumn{3}{c}{$+45$ ($+2.9\%$)} &
       \multicolumn{3}{c}{$+3$ ($+4.5\%$)} \\
\bottomrule
\multicolumn{10}{l}{\footnotesize All $p<0.001$ (paired Wilcoxon, 30 seeds). Lift = Ours $-$ Opt-only.
  TRV values $<\$100$K shown as rounded integers.}\\
\multicolumn{10}{l}{\footnotesize $p$-value reflects simulation stability across seeds, not cross-dataset
  generalization.}\\
\end{tabular}
\end{table*}

\smallskip
\noindent\textbf{S1 interpretation.}
In IT infrastructure decommission, the framework's TRV advantage over
Opt-only (\$73K, $+13.7\%$) comes largely from two effects: the keyword
extractor correctly flags damaged assets early, preventing wasted disassembly
effort; and the adaptive inspection policy skips \$32K in inspection cost that
the inspect-everything baselines must pay regardless. Full demand coverage
(DCR~100\%) reflects that with strong text signal, the optimizer confidently
processes every asset to the right recovery level.

\smallskip
\noindent\textbf{S2 interpretation.}
Aviation MRO is the most challenging scenario. The TRV lift is positive but
modest (\$45K, $+2.9\%$), and XGBoost outperforms Ours in absolute TRV
(\$1,678K vs.\ \$1,610K). This is consistent with the weak text calibration
($r=0.32$): when the text signal is marginal, the inspection cost savings from
the adaptive policy do not fully compensate for the occasional missed recovery
that a fully-inspecting baseline would catch. The framework saves \$10K in
inspection cost over Opt-only (\$34K vs.\ \$24K), confirming that it is
correctly routing low-confidence assets to full inspection. In a regulated
environment, this advisory behavior—high-confidence skip plus full inspection
for uncertain cases—is the appropriate operating mode.

\smallskip
\noindent\textbf{S3 interpretation.}
Consumer electronics returns show the triage pattern. Sem.-only (threshold
routing without optimization) produces negative TRV because inspection costs
exceed recovery value for the assets it selects; Ours avoids this by
combining the text signal with the greedy optimizer, achieving positive lift
(\$3K, $+4.5\%$) over Opt-only. At 1000 units per batch with \$20 inspection
costs, the per-unit economics make even small categorical routing improvements
meaningful at volume.

\smallskip
\noindent\textbf{Statistical note.}
With a deterministic extractor, randomness enters only through Monte Carlo yield
sampling and seed-varied asset populations. The paired Wilcoxon $p<0.001$ across
30~seeds measures stability of these effects, not generalization to held-out
datasets or to a different extractor. Every value in Table~\ref{tab:main}
reproduces by running \texttt{python scripts/run\_summary.py --seeds 0-29} in
the released RLDB code.

\subsection{Extractor Calibration}
\label{sec:calibration}

Table~\ref{tab:calibration} reports Pearson $r$ between extractor output $\phi$
and realized yield across 30~seeds (pooled; 15,000 asset-evaluations for S1/S2,
30,000 for S3). Two extractors are compared: the keyword classifier used in the
main simulation, and a deterministic oracle-text reader that correctly identifies
the described condition from the note text and maps it to a recoverability
estimate via a table fixed before extraction. The oracle reader represents the
upper bound achievable by any deterministic classifier on the same synthetic
benchmark, given the irreducible noise introduced by the decoupled note-generation
process ($p_\text{omit}=0.15$, $p_\text{mislabel}=0.25$).

\begin{table}[t]
\caption{Extractor Calibration on the Synthetic Benchmark\\
  ($r$ = Pearson correlation between $\phi$ and true yield, pooled over 30 seeds)}
\label{tab:calibration}
\centering
\begin{tabular}{lcrr}
\toprule
Scenario & $N$ & Keyword & Oracle-text \\
         &     & (lower bound) & (upper bound) \\
\midrule
S1: IT Infrastructure    & 15,000 & 0.47 & 0.73 \\
S2: Aviation MRO         & 15,000 & 0.32 & 0.54 \\
S3: Consumer Electronics & 30,000 & 0.36 & 0.74 \\
\bottomrule
\multicolumn{4}{l}{\footnotesize Both extractors are deterministic; neither is a
  language model.}\\
\multicolumn{4}{l}{\footnotesize Note generation is decoupled from ground-truth
  yield (see text).}\\
\end{tabular}
\end{table}

Several features of Table~\ref{tab:calibration} deserve attention. First, the
gap between lower and upper bounds is substantial (0.26--0.38 Pearson units),
indicating that a better reader of the same text would meaningfully improve
framework performance. This is the extractor-agnostic property in action: the
decision architecture is the same; only the extraction quality changes. Second,
the fallback rate—the fraction of notes for which no signal fires and $\phi$
defaults to~1.0—is 28\% for S1, 43\% for S2, and 55\% for S3, explaining the
ordering: more colloquial and abbreviated text generates more uninformative notes.
Third, the bounds differ across scenarios, ruling out any concern that the numbers
are the result of a shared evaluation artifact.

\smallskip
\noindent\textbf{S1 note.} Even with decoupled generation, the keyword extractor's
$r=0.47$ is better than chance, confirming that the stylized vocabulary captures
real signal. The ceiling ($r=0.73$) is the theoretical maximum on this benchmark
given the imposed note noise; organic technician notes would have different
(likely higher) noise characteristics and would require recalibration.

\smallskip
\noindent\textbf{S2 note.} Standardized ATA codes create a flat distribution
of signal words that the keyword classifier cannot order reliably. The ceiling
is 0.54 rather than 0.73 because even a perfect reader cannot distinguish
conditions that the coding system collapses into the same vocabulary.

\subsection{Error Analysis}

We identify errors as asset-level decisions where the framework's realized
recovery is more than \$50 below what the optimization-only baseline would
have achieved for the same asset. The dominant error type is overestimation:
the extractor assigns a higher $\phi$ than warranted, causing the system to
skip inspection on assets that should have been tested. Error rates are
approximately 9\% in S1 (45/500 assets), 24\% in S2 (118/500), and near-zero
in S3. Total loss from errors is \$31K (S1), \$70K (S2), and negligible (S3),
well below the semantic lift in all scenarios, confirming net positive value.
The S2 error rate reflects the standardized-code calibration ceiling and is
the primary reason aviation applications must keep a human in the loop.

\subsection{Case Studies}

\noindent\textbf{Case 1 --- Correct skip (S1).} A note reads
``\textit{Routine decommission. All components seated properly. No corrosion.}''
The extractor returns $\phi=0.95$, $\sigma=0.95$. High confidence triggers an
inspection skip; the asset proceeds directly to $L_2$ disassembly, recovering
\$4,200 in components while saving \$200 in inspection cost.

\noindent\textbf{Case 2 --- Correct flag (S1).} A note reads
``\textit{PSU failure. Visible burn marks on mainboard near power connector J12.
CPU smells burnt.}'' The extractor returns $\phi=0.05$, $\sigma=0.95$.
Despite high confidence, the low context factor routes the asset to scrap,
avoiding \$70 in wasted disassembly. Confidence and condition are distinct axes:
the model is certain the unit is damaged.

\noindent\textbf{Case 3 --- Overestimation failure (S2).} An SDR reads
``\textit{Oil on bottom of \#1 engine fan cowl. Replaced filter and packing.}''
The keyword extractor detects ``replaced'' and emits $\phi=0.98$, $\sigma=0.85$.
The true yield factor is 0.59—the repair addressed a symptom, not the root cause.
This overestimation costs \$8,669 against the full-inspection baseline and
illustrates exactly why keyword matching cannot safely replace physical inspection
in safety-critical MRO: a simple word cannot infer whether a repair resolved the
underlying defect.

\noindent\textbf{Case 4 --- Categorical triage (S3).} A review reads
``\textit{It just stopped working after a week. Screen never turned on.}''
Continuous yield prediction on this text is near-zero, but the extractor
correctly bins the unit as ``likely dead'' and routes it to component harvesting
rather than whole-unit refurbishment. The value is categorical routing at
volume, not a precise yield number.

\subsection{Sensitivity Analysis}

\noindent\textbf{Capacity constraints.}
Under tight capacity (50\% of baseline), S1 TRV drops 43\% and S2 drops 23\%,
confirming selective harvesting: when capacity binds, the engine concentrates
effort on the highest-value assets and scraps the rest. This responsiveness to
capacity is a direct benefit of the value-ranked greedy allocator.

\noindent\textbf{Confidence thresholds.}
Sweeping $\tau_h$ from 0.5 to 0.95 and $\tau_l$ from 0.2 to 0.75: S1 and S2
favor $\tau_h=0.5$, $\tau_l=0.25$ (aggressive skip); S3 favors a narrower gap
($\tau_h=0.5$, $\tau_l=0.45$) because colloquial text rarely provides very high
confidence. Even moderate extractor confidence is enough to skip inspection
profitably because the expected cost of a wrong skip (\$200--700) is often less
than the certain cost of universal inspection (\$120--200 per asset).

\noindent\textbf{$\kappa$ sensitivity.}
Table~\ref{tab:kappa} reports TRV lift across $\kappa \in \{5, 10, 20, 40, 60, 80\}$
on S1. The near-flat profile (spread: 0.86 percentage points) confirms the
analysis in Section~\ref{sec:framework-prior}: framework value derives from
inspection targeting, not from prior concentration. We retain $\kappa=10$ as
a conservative default.

\begin{table}[t]
\caption{Effect of $\kappa$ on S1 TRV Lift (30 seeds)}
\label{tab:kappa}
\centering
\begin{tabular}{lcccccc}
\toprule
$\kappa$ & 5 & 10 & 20 & 40 & 60 & 80 \\
\midrule
TRV lift (\%) & 14.0 & 13.7 & 14.1 & 14.2 & 14.6 & 13.9 \\
\bottomrule
\multicolumn{7}{l}{\footnotesize Spread 0.86\,pp. Flatness is structural: decisions use the
  Beta mean, which is}\\
\multicolumn{7}{l}{\footnotesize $\kappa$-invariant.}\\
\end{tabular}
\end{table}

%---------------------------------------------------------------------------
\section{Discussion}

\subsection{Practical Implications}

The framework requires minimal infrastructure: a greedy allocator with MC
sampling (no specialized solver), and any extractor satisfying a simple
interface. Adapting to a new industry means updating the extraction vocabulary
and recalibrating the $\phi$-to-yield mapping—not retraining a model. This
makes the approach accessible to operators without ML engineering resources.

The strongest case for deployment is where text is already generated as part of
normal operations (IT decommissions, MRO work orders) and inspection is the
binding constraint. In these settings, even the conservative keyword extractor
yields meaningful TRV gains while substantially cutting inspection cost.

\subsection{Limitations}

Several limitations bound the scope of these findings. First, all simulation
data is synthetic. Calibration parameters come from public sources, but the
generators produce synthetic notes and condition values; organic text from real
operations may have different noise characteristics that require extractor
recalibration. Second, the S2 error rate (approx.\ 24\%) confirms that the
framework must not operate autonomously in safety-critical domains: any
deployment should preserve mandated physical inspection as the final decision
authority, with the semantic signal serving only to direct triage and prioritize
inspection effort. Third, the optimization is single-period; multi-period
dynamics with inventory carry-over remain open. Fourth, statistical significance
($p<0.001$, 30~seeds) measures simulation stability across populations, not
generalization to new datasets or text distributions.

\subsection{RLDB Benchmark}

The Reverse Logistics Decision Benchmark (RLDB) released with this paper
contains: scenario configurations for all three settings; stochastic asset
generators with decoupled note-condition generation; calibration scripts that
produce all reported tables without manual intervention; and implementations of
all five baselines. Running \texttt{python scripts/run\_summary.py --seeds 0-29}
reproduces Table~\ref{tab:main}; \texttt{run\_calibration.py} reproduces
Table~\ref{tab:calibration}; \texttt{run\_kappa\_sweep.py} reproduces
Table~\ref{tab:kappa}. Repository: \href{https://github.com/shuan63/reverse_logistics}{github.com/shuan63/reverse\_logistics}.

%---------------------------------------------------------------------------
\section{Conclusion}

Integrating semantic signals from unstructured return text into stochastic
disposition optimization is both feasible and beneficial across a range of
reverse logistics settings. The key mechanism is an adaptive inspection policy
that uses extractor confidence to determine inspection depth per asset, replacing
uniform inspect-everything with targeted physical testing only where text evidence
is weak. Even a conservative keyword extractor—our deliberate lower bound—
produces positive TRV lifts of 2.9--13.7\% over an inspect-all baseline while
saving meaningful inspection cost.

The honest picture is that the benefit is concentrated where text is genuinely
informative. In IT infrastructure with descriptive technician notes, the framework
is clearly worth deploying. In aviation MRO, where standardized codes limit text
signal, it is a useful advisory tool but not a replacement for mandated physical
inspection. In consumer electronics, where per-unit inspection is often
economically infeasible regardless, the triage value—routing returns to the
right quality tier at volume—is the primary contribution.

The framework is extractor-agnostic by design. Every result presented here is a
lower bound; a stronger extractor plugs directly into the same architecture and
improves performance proportionally to its calibration gain. That direct link
between text quality and decision quality, confirmed across three structurally
diverse scenarios, is the central empirical finding of this work.

%---------------------------------------------------------------------------
\bibliographystyle{IEEEtran}
\bibliography{refs}

\end{document}
